\chapter{Un vicolo cieco istruttivo}
\label{cap:vicolo_cieco_corridoi}

Il capitolo precedente si è chiuso con un'idea che sembrava promettente: calcolare, da un modello del terreno, i corridoi a minor costo energetico seguiti dagli animali, usando un modello fisico solido come quello di Pontzer, e usarli per individuare i punti di passaggio più probabili, dove i cacciatori preistorici avrebbero potuto attendere le proprie prede. Questo capitolo racconta perché quell'idea, messa alla prova, si sia rivelata meno utile del previsto, e cosa quel fallimento abbia insegnato.

\section{Il problema nascosto nei numeri}

L'idea di un corridoio di spostamento presuppone qualcosa di preciso: che l'animale in questione si muova su distanze relativamente lunghe, magari stagionalmente, passando da un'area all'altra lungo un tragitto ricorrente. È un'immagine intuitiva, rafforzata dal fatto che oggi si parla spesso, a proposito di grandi predatori come il lupo o l'orso, di territori vastissimi percorsi nell'arco di una stagione.

I dati raccolti nel quinto capitolo raccontano, però, una storia diversa per le specie che più interessano il problema di questo libro: stambecco e camoscio, le prede dominanti nei siti di caccia più tipici delle Alpi venete preistoriche, come mostrato nel dettaglio dai dati di Riparo Dalmeri e Riparo Soman nel sesto capitolo. Come si è visto, entrambe le specie occupano territori piccoli, nell'ordine di uno o pochi chilometri quadrati per il camoscio (spesso meno di un chilometro quadrato), di poche decine per lo stambecco, con una fedeltà al proprio territorio che porta gli stessi individui a tornare, anno dopo anno, negli stessi luoghi.

Questo dato smonta, alla radice, l'idea di un lungo corridoio di spostamento da modellare come un percorso a minor costo energetico tra due punti lontani. Un camoscio non percorre un tragitto di molti chilometri tra un'area di svernamento e un'area estiva: resta, sostanzialmente, nello stesso territorio ristretto tutto l'anno, con al più uno spostamento verticale limitato legato alla coltre nevosa invernale. Calcolare un corridoio a costo minimo tra due punti lontani, per un animale che in realtà non compie quel tragitto, significa costruire un modello elegante che risponde a una domanda sbagliata.

\section{Un problema che la ricerca ecologica riconosce esplicitamente}

Questo non è solo un problema specifico del ragionamento di questo libro: è un limite riconosciuto nella letteratura scientifica dedicata alla modellazione dei corridoi faunistici. Una rassegna recente sulla pianificazione della connettività ecologica osserva esplicitamente che gli animali di grande taglia e ampio raggio d'azione, spesso chiamati specie ``flagship'' o ``ombrello'', sono quelli storicamente più studiati nella ricerca sui corridoi, proprio perché si presume forniscano un beneficio di connettività utile anche per altre specie \citep{bohnett2024}. La stessa rassegna avverte che una pianificazione di corridoio efficace per una specie può non esserlo affatto per un'altra, proprio a causa delle differenze di idoneità dell'habitat, di distanza di dispersione e di estensione del territorio.

Questo significa che il modello di corridoio a costo minimo, per come è tipicamente costruito e validato in letteratura, è tarato su un tipo di animale (grande, mobile, con home range esteso) che è quasi l'opposto di camoscio e stambecco. Applicarlo a queste specie non è solo empiricamente poco utile, come mostrano i dati di home range discussi nel quinto capitolo: è metodologicamente improprio, un caso di applicazione di uno strumento fuori dal contesto per cui è stato pensato e validato.

\section{L'eccezione che complica, non che salva}

Il cervo, come si è visto nel quinto capitolo, si comporta diversamente: alcuni individui migrano stagionalmente, spostandosi in modo più marcato tra quote diverse, con home range che nella stagione calda possono superare abbondantemente quelli invernali. Per questa specie, un modello di corridoio a costo minimo avrebbe più senso ecologico, ed è coerente con il tipo di specie su cui la modellazione dei corridoi funziona meglio secondo la letteratura appena citata.

Ma questo non salva l'idea originaria, la complica soltanto: significa che un modello unico di corridoio, applicato indistintamente a tutta la fauna cacciata, mescola comportamenti radicalmente diversi (specie residenti a raggio ridotto, specie migratrici a raggio ampio) sotto un'unica formula, con il rischio concreto di produrre un risultato che non descrive bene nessuna delle due situazioni. E i dati archeozoologici discussi nel sesto capitolo mostrano che, nei siti prealpini veneti meglio documentati, sono proprio stambecco e camoscio le specie dominanti, non il cervo: a Riparo Dalmeri lo stambecco supera da solo il 90\% dei resti, a Riparo Soman il camoscio supera il 40\%. Un modello di corridoio, anche se teoricamente più adatto al cervo, finirebbe per descrivere bene la specie meno rilevante per la maggior parte dei siti conosciuti.

\section{Un secondo tentativo, altrettanto instabile}

Di fronte a questo primo intoppo, un secondo tentativo si è concentrato su un ingrediente diverso: invece di modellare il movimento, modellare la presenza di risorse necessarie alla sopravvivenza, in particolare acqua e vegetazione, per individuare le aree più favorevoli alla fauna. È un'idea mutuata da contesti geografici diversi, dove funziona bene: nella Región de Atacama discussa nel primo capitolo, per esempio, gli animali selvatici si concentrano attorno a risorgive e zone umide che restano nella stessa posizione per millenni, offrendo un punto di riferimento stabile e ricostruibile.

Questo secondo tentativo si è arenato quasi subito su un ostacolo pratico, non teorico. La vegetazione alpina di dodicimila o quattordicimila anni fa non è ricostruibile con la precisione necessaria a un modello di questo tipo: si sa, in linea generale, come è cambiato nel tempo il limite del bosco (il quarto capitolo lo ha già descritto, con le sue soglie approssimative dichiarate), ma non si sa, punto per punto, che tipo di copertura vegetale ci fosse in un dato luogo in un dato momento, informazione che richiederebbe dati palinologici puntuali che semplicemente non esistono per la maggior parte del territorio, al di fuori dei pochi punti, come il bacino di Palughetto sul Cansiglio, dove una sequenza del genere è stata effettivamente ricostruita.

Anche il criterio dell'acqua, in questo contesto, si è rivelato meno utile di quanto sembrasse. A differenza delle risorgive andine, stabili nel tempo, il reticolo idrografico alpino è stato pesantemente rimodellato nel corso dei millenni da frane, alluvioni ed erosione: la posizione esatta di un piccolo corso d'acqua o di una sorgente diecimila anni fa non è affatto garantita coincidere con quella attuale.

\section{Cosa resta di questi due tentativi}

Raccontare questi due tentativi falliti non è un esercizio di autocritica fine a se stesso. Serve a mostrare un punto metodologico che vale più di qualunque singolo risultato, e che trova ora conferma sia nei dati zoologici specifici discussi nel quinto capitolo sia nella letteratura generale sulla modellazione dei corridoi faunistici: un modello elegante, costruito su basi fisiche corrette (come la relazione tra costo energetico e pendenza descritta dal modello di Pontzer), può comunque fallire se applicato a un problema che non corrisponde alla realtà del comportamento animale che si vuole descrivere, o se applicato a un tipo di specie diverso da quello per cui il metodo è stato pensato. La fisica del movimento era giusta; l'idea che stambecco e camoscio percorressero lunghi corridoi era sbagliata, e lo era per una ragione riconosciuta anche al di fuori del caso specifico di questo libro. La logica dell'acqua come risorsa chiave era corretta; l'idea di poter ricostruire con precisione dove si trovasse quell'acqua migliaia di anni fa era eccessivamente ottimistica.

Il criterio che serve, allora, non è un percorso, ma qualcos'altro: un modo per riconoscere, sul terreno di oggi, i luoghi che più probabilmente corrispondevano ai territori ristretti e fedeli di stambecco e camoscio, senza pretendere di ricostruire un tragitto che quegli animali, con ogni probabilità, non compivano affatto. Il prossimo capitolo mostra come qualcuno, lavorando su un territorio alpino diverso ma per molti versi comparabile, abbia affrontato esattamente questo problema, con un metodo che ha il pregio di essere fondato su dati concreti raccolti sul campo, non su un modello teorico costruito a tavolino.
